\section{A Hope to Renew Each Year}

An annual return to heaven can gather the scattered desires of a life. During the year many immediate concerns rightly claim attention, while others acquire an urgency they do not deserve. A persistent worry can occupy our attention so completely that the promised end seems remote. To consider heaven again is to recover proportion. The question is no longer merely whether the year has gone as we wished. It is whether we have allowed its goods and trials to draw us toward the communion for which we were made, and whether we are ready to receive the grace still offered.

The Church's remembrance of the saints and the faithful departed gives this reflection a human company. The saints show what grace can accomplish in persons who once faced weakness, uncertainty, and the demands of their own time. The dead whom we commend to God remind us that the pilgrimage has a real end. We can thank him for their gifts without making ourselves judges of their final state. We can pray for their purification and eternal rest, and renew our own willingness to be purified. The communion for which we hope is already teaching us to love across the boundary of death.

Francis de Sales brings his meditation to a definite resolution because a desire left wholly indefinite can be displaced by the next attraction. The return to heaven may disclose one attachment that repeatedly rules us or one neglected good through which charity could grow. We need not produce an elaborate account of the coming year. A faithful response can begin with a concrete consent that is possible today. Keeping that resolution gives the desire a place in the decisions of the coming day. The Lord forms persons through the ordinary matter of their lives, and a repeated fidelity can become the place where a divided heart learns to love.

Paul's desire to depart and be with Christ stands beside his readiness to remain for the good of the Philippians (Phil 1:20--26). His longing does not make their need an interruption of his true vocation. Christ is his life, whether that life is spent in further service or brought to the fulfillment he desires. This is a mature form of heavenly hope. We may long to be with the Lord while receiving the present day as the day in which he asks us to love. The timing of our summons belongs to him; the charity of our response belongs to the grace he gives now.

This hope can also remain when spiritual feeling is quiet. We may believe sincerely and still feel more apprehension than longing at the thought of death. Affection for those who depend upon us, bodily fear, and sorrow over unfinished responsibilities can coexist with trust. We can place these movements before God without pretending they have already been resolved. The desire to desire heaven is itself something we can bring to prayer. Christ's promise is firmer than the changing strength of our response, and his patient work can deepen our readiness in ways that are hidden even from us.

God promises to make love whole. The blessed will know him directly, love him without sin, rejoice with his saints, and share the risen life of Christ in a creation made new. We await the fulfillment of these promises while receiving the grace to persevere today. We are learning to receive the one who has already given himself for us and who desires that we should be with him.

The time still given to us can be spent answering Christ's promise: returning to him after sin, receiving his gifts, and serving those entrusted to us. We cannot know when that time will end. We can ask to remain faithful within it and to desire more fully the life he offers. The God who already draws us into friendship is preparing us to rest in him, see him, love him, and praise him without end.

